\documentclass[12pt]{article}
\usepackage{ficha}

\begin{document}

\ficha{Operações com Números Decimais}

\instrucoes{Resolve sem calculadora, apresentando os cálculos (podes usar o algoritmo em coluna).}

\secao{Adição e Subtração}{Calcula.}
\begin{exercicios}
  \item $3{,}4 + 2{,}5$
  \item $12{,}75 + 4{,}3$
  \item $0{,}8 + 0{,}25$
  \item $15 + 3{,}06$
  \item $9{,}6 - 4{,}2$
  \item $10 - 3{,}45$
  \item $7{,}05 - 2{,}8$
  \item $25{,}3 - 18{,}75$
  \item $4{,}5 + 2{,}35 - 1{,}9$
  \item $100 - 0{,}01$
\end{exercicios}

\secao{Multiplicação e Divisão por 10, 100, 1000 e por 0,1, 0,01, 0,001}{Calcula.}
\begin{exercicios}
  \item $3{,}45 \times 10$
  \item $0{,}067 \times 100$
  \item $2{,}5 \times 1\,000$
  \item $48{,}2 : 10$
  \item $7 : 100$
  \item $356 : 1\,000$
  \item $64 \times 0{,}1$
  \item $5{,}3 \times 0{,}01$
  \item $12 : 0{,}1$
  \item $0{,}9 : 0{,}01$
  \item $1{,}2 \times 0{,}001$
  \item Completa: $4{,}7 \times \square = 470$ \quad e \quad $82 : \square = 0{,}082$.
\end{exercicios}

\secao{Multiplicação}{Calcula.}
\begin{exercicios}
  \item $2{,}3 \times 4$
  \item $0{,}6 \times 7$
  \item $1{,}5 \times 1{,}2$
  \item $3{,}25 \times 4$
  \item $0{,}4 \times 0{,}3$
  \item $2{,}05 \times 3{,}1$
  \item $12{,}5 \times 0{,}8$
  \item $0{,}07 \times 0{,}5$
  \item $(1{,}5)^2$
  \item $4{,}8 \times 2{,}5$
\end{exercicios}

\secao{Divisão}{Calcula.}
\begin{exercicios}
  \item $8{,}4 : 4$
  \item $9{,}6 : 3$
  \item $7 : 2$
  \item $3 : 4$
  \item $1{,}2 : 0{,}3$
  \item $4{,}5 : 0{,}5$
  \item $0{,}72 : 0{,}08$
  \item $15 : 0{,}25$
  \item $6{,}3 : 1{,}4$
  \item $2{,}88 : 1{,}2$
\end{exercicios}

\secao{Expressões e Problemas}{Calcula ou resolve.}
\begin{exercicios}
  \item $2{,}5 + 1{,}5 \times 4$
  \item $(2{,}5 + 1{,}5) \times 4$
  \item $10 - 0{,}2 \times 15$
  \item $3{,}6 : 1{,}2 + 0{,}5 \times 0{,}4$
  \item Um caderno custa 2,35 euros. Quanto custam 6 cadernos?
  \item Um rolo de 12,6\,m de fio foi dividido em 9 partes iguais. Quanto mede cada parte?
  \item A Maria comprou 1,5\,kg de maçãs a 1,80 euros/kg e 0,75\,kg de queijo a 12 euros/kg. Quanto pagou?
  \item Um carro gasta 6,4 litros por cada 100\,km. Quantos litros gasta em 350\,km?
\end{exercicios}

\secao{Desafios}{Resolve os problemas.}
\begin{exercicios}
  \item O produto de dois números é 3,6. Se um deles for 0,9, qual é o outro?
  \item Quantas garrafas de 0,33 litros se enchem com 9,9 litros de sumo?
  \item Uma sala retangular mede 4,5\,m por 3,2\,m. O chão custa 18,50 euros por m$^2$. Quanto custa revestir o chão?
  \item O Tiago tinha 20 euros. Comprou 3 sandes a 2,75 euros e 2 sumos a 1,40 euros. Quanto recebeu de troco?
  \item Multiplicar um número por 0,5 é o mesmo que dividi-lo por quanto? E multiplicar por 0,25? Justifica.
  \item Dividir por 0,1 aumenta ou diminui o número? Explica com um exemplo.
  \item Encontra dois números decimais cuja soma seja 5 e cujo produto seja 6,24.
  \item Um atleta corre 400\,m em 52,8 segundos. Quanto tempo demora, em média, a correr 100\,m?
\end{exercicios}

\end{document}
